\documentclass[11pt]{article}

\usepackage[margin=1in]{geometry}
\usepackage{amsmath,amssymb,amsthm}
\usepackage{bm}
\usepackage{booktabs}
\usepackage{mathtools}
\usepackage[colorlinks=true,linkcolor=blue,citecolor=blue]{hyperref}

\newtheorem{theorem}{Theorem}
\newtheorem{lemma}{Lemma}
\newtheorem{proposition}{Proposition}
\theoremstyle{remark}
\newtheorem{remark}{Remark}

\newcommand{\R}{\mathbb{R}}
\newcommand{\E}{\mathbb{E}}
\newcommand{\tr}{\operatorname{tr}}
\newcommand{\vecop}{\operatorname{vec}}
\newcommand{\Wqk}{W_{qk}}
\newcommand{\Wov}{W_{ov}}
\newcommand{\Wq}{W_{q}}
\newcommand{\Wk}{W_{k}}
\newcommand{\Wo}{W_{o}}
\newcommand{\Wv}{W_{v}}
\newcommand{\softmax}{\operatorname{softmax}}
\newcommand{\Hess}{\nabla^2}
\DeclareMathOperator{\rank}{rank}

\title{\bfseries Why the Combined $(\Wqk,\Wov)$ Parametrization Has a More\\
Favorable Loss Landscape than the Factored $(\Wq,\Wk,\Wo,\Wv)$ One}
\author{ICL project --- folder 7 (\texttt{unique\_Wqk\_Wok\_test})}
\date{\today}

\begin{document}
\maketitle

\begin{abstract}
A disentangled attention layer can be parametrized in two functionally
equivalent ways: \emph{combined}, directly by the circuit matrices
$\Wqk=\Wq^{\!\top}\Wk$ (scores) and $\Wov=\Wo\Wv$ (values); or \emph{factored},
by the four matrices $\Wq,\Wk,\Wo,\Wv$. We show that the factored
loss is the combined loss \emph{pulled back through a bilinear map}, and that
this pullback can only \emph{add} curvature pathologies, never remove them.
Precisely: the factored Hessian equals a positive‑semidefinite Gauss--Newton
pullback of the combined Hessian \emph{plus} an indefinite term whose
eigenvalues are $\pm\sigma_\ell(\nabla_M\widetilde{L})$ --- the singular values of
the reduced gradient. Consequently the factored landscape (i)~carries a
$r^2{+}r'^2$‑dimensional flat \emph{gauge} null space per layer, and (ii)~develops
strictly \emph{negative} curvature (saddles) at every non‑stationary point,
exactly where gradient descent operates. We confirm this empirically on the
trained models: the combined Hessian is near‑isotropic (condition number
$\kappa\!\approx\!1.2$) and stays positive semidefinite off the optimum, whereas
the factored Hessian is $\kappa\!\approx\!10^4$ and acquires a negative
eigenvalue $\lambda_{\min}\!\approx\!-0.74$ off the optimum. The combined
parametrization is therefore the more convex, better‑conditioned, and
more gradient‑descent‑favorable coordinate system.
\end{abstract}

\section{Setup: one attention layer, two coordinate systems}\label{sec:setup}

Let $X\in\R^{T\times d}$ be the (token $\oplus$ position) representation that
enters an attention layer, with $T$ positions and width $d$. The disentangled
attention layer used throughout this project computes, with causal mask
$\mathcal M$ and head width $m$,
\begin{equation}\label{eq:layer}
S \;=\; \tfrac{1}{\sqrt m}\, X\,\Wqk\,X^{\!\top} + \mathcal M,
\qquad
\text{attn} \;=\; \softmax(S)\,\bigl(X\,\Wov^{\!\top}\bigr),
\qquad
\text{out} = [\,X \,\|\, \text{attn}\,].
\end{equation}
The \textbf{combined parametrization} treats the two circuit matrices as the
free variables,
\[
\boxed{\;\Wqk\in\R^{d\times d}\quad(\text{score circuit}),\qquad
        \Wov\in\R^{m\times d}\quad(\text{value circuit}).\;}
\]
The \textbf{factored (separate) parametrization} writes each circuit as a product
of two matrices,
\begin{equation}\label{eq:factor}
\Wqk=\Wq^{\!\top}\Wk,\quad \Wq,\Wk\in\R^{r\times d},
\qquad\qquad
\Wov=\Wo\,\Wv,\quad \Wo\in\R^{m\times r'},\ \Wv\in\R^{r'\times d},
\end{equation}
so that the score becomes $S=\tfrac{1}{\sqrt m}X\Wq^{\!\top}\Wk X^{\!\top}+\mathcal M$
and the value path becomes $\softmax(S)\,(X\Wv^{\!\top}\Wo^{\!\top})$. Both
parametrizations realize \emph{exactly the same function} as long as the inner
ranks are large enough ($r\ge\rank\Wqk$, $r'\ge\rank\Wov$); the standard
attention head is the special case $r=r'=m$.

Stack the combined variables into $M=(\Wqk,\Wov)$ and the factored variables into
$\theta=(\Wq,\Wk,\Wo,\Wv)$. The map relating them,
\begin{equation}\label{eq:phi}
\varphi(\theta)=M,\qquad
\varphi_{qk}(\Wq,\Wk)=\Wq^{\!\top}\Wk,\quad
\varphi_{ov}(\Wo,\Wv)=\Wo\Wv,
\end{equation}
is \textbf{bilinear} (linear in each factor separately). This single algebraic
fact is the source of everything below.

Let $\widetilde L(M)=\E[\text{cross-entropy}]$ be the (suffix-$100$) training loss
as a function of the circuits. The two landscapes are
\begin{equation}\label{eq:losses}
L_{\mathrm{comb}}(M)=\widetilde L(M),
\qquad\qquad
L_{\mathrm{sep}}(\theta)=\widetilde L\bigl(\varphi(\theta)\bigr).
\end{equation}
They describe the \emph{same} surface, but in different coordinates. Gradient
descent moves along the coordinates it is given, so the geometry that matters is
the geometry of \eqref{eq:losses} in each coordinate system.

\section{The landscape is parametrization-dependent: the Hessian chain rule}

Curvature is not coordinate-free. For the composition
$L_{\mathrm{sep}}=\widetilde L\circ\varphi$, the exact first and second
derivatives are
\begin{align}
\nabla_\theta L_{\mathrm{sep}}
   &= J_\varphi^{\!\top}\,\nabla_M\widetilde L, \label{eq:grad}\\[2pt]
\Hess_\theta L_{\mathrm{sep}}
   &= \underbrace{J_\varphi^{\!\top}\,\bigl(\Hess_M\widetilde L\bigr)\,J_\varphi}
        _{\textstyle \text{Gauss--Newton pullback }=:\,P}
     \;+\;
      \underbrace{\sum_{a}\bigl(\nabla_M\widetilde L\bigr)_a\,\Hess_\theta\varphi_a}
        _{\textstyle \text{residual curvature }=:\,R}, \label{eq:hess}
\end{align}
where $J_\varphi=\partial\,\vecop(M)/\partial\,\vecop(\theta)$ is the Jacobian of
\eqref{eq:phi} and $a$ indexes the entries of $M$. The combined Hessian is just
$\Hess_M\widetilde L$. The two extra structures in \eqref{eq:hess} --- the
\emph{pullback} $P$ and the \emph{residual} $R$ --- are precisely what the
factored coordinates inject. We analyze each.

\section{The curvature of a bilinear factorization}

The residual term $R$ is governed entirely by the bilinearity of $\varphi$.

\begin{lemma}[Curvature of $A\mapsto A^{\!\top}B$]\label{lem:bilin}
Fix $G\in\R^{d\times d}$ and consider $\psi(A,B)=\langle G,\,A^{\!\top}B\rangle
=\tr(G^{\!\top}A^{\!\top}B)$ with $A,B\in\R^{r\times d}$. Then $\psi$ is a quadratic
form whose Hessian in $\theta=(\vecop A,\vecop B)$ is the block matrix
\begin{equation}\label{eq:Rblock}
\Hess_\theta\psi \;=\;
\begin{bmatrix} 0 & I_r\otimes G\\[2pt] I_r\otimes G^{\!\top} & 0\end{bmatrix},
\end{equation}
i.e.\ the $A$--$A$ and $B$--$B$ blocks vanish and only the cross block survives.
Its eigenvalues are
\[
\boxed{\;\pm\,\sigma_\ell(G),\quad \ell=1,\dots,\rank(G),\ \text{each of multiplicity }r,\;}
\]
together with zeros. In particular $\Hess_\theta\psi$ is \emph{indefinite}
whenever $G\neq 0$, with smallest eigenvalue $-\sigma_{\max}(G)$.
\end{lemma}

\begin{proof}
Writing $\psi=\sum_{i,j}G_{ij}\sum_{k}A_{ki}B_{kj}$, the first derivatives are
$\partial\psi/\partial A = BG^{\!\top}$ and $\partial\psi/\partial B = AG$, and the
only nonzero second derivatives are the cross terms
$\partial^2\psi/\partial A_{ki}\,\partial B_{lj}=\delta_{kl}\,G_{ij}$, which
assemble into \eqref{eq:Rblock}. A block matrix
$\bigl[\begin{smallmatrix}0&C\\ C^{\!\top}&0\end{smallmatrix}\bigr]$ has
eigenvalues $\pm\sigma(C)$; here $C=I_r\otimes G$ has singular values
$\{\sigma_\ell(G)\}$ each repeated $r$ times.
\end{proof}

Because $\widetilde L$ depends on $\Wq,\Wk$ only through $\Wqk=\Wq^{\!\top}\Wk$, the
QK part of the residual in \eqref{eq:hess} is exactly $\Hess_\theta\psi$ of
Lemma~\ref{lem:bilin} with the \textbf{reduced gradient}
\begin{equation}\label{eq:redgrad}
G \;=\; \nabla_{\Wqk}\widetilde L \;\in\;\R^{d\times d}.
\end{equation}
The OV part is identical with $\varphi_{ov}(\Wo,\Wv)=\Wo\Wv$ and
$G'=\nabla_{\Wov}\widetilde L$. Hence

\begin{equation}\label{eq:Rspec}
\text{spectrum}(R)\;=\;\{\pm\sigma_\ell(\nabla_{\Wqk}\widetilde L)\}
                    \,\cup\,\{\pm\sigma_\ell(\nabla_{\Wov}\widetilde L)\}\,\cup\,\{0\}.
\end{equation}

\paragraph{Reading \eqref{eq:Rspec}.}
The factored coordinates contribute a curvature term whose negative eigenvalues
have magnitude equal to the \emph{singular values of the reduced gradient}. This
term \emph{vanishes at stationary points} of $\widetilde L$ (where
$\nabla_M\widetilde L=0$) and is \emph{strictly indefinite everywhere else} ---
that is, at every point gradient descent actually visits on its way to a
solution.

\section{Gauge redundancy: flat valleys with no analogue in the combined map}

The factorization \eqref{eq:factor} is not unique. For any invertible
$P\in GL(r)$,
\begin{equation}\label{eq:gauge}
(\Wq,\Wk)\ \longmapsto\ (P^{-\top}\Wq,\ P\Wk)
\qquad\Longrightarrow\qquad
\Wq^{\!\top}\Wk \ \text{unchanged},
\end{equation}
and analogously $(\Wo,\Wv)\mapsto(\Wo Q^{-1},Q\Wv)$ for $Q\in GL(r')$. Along these
\textbf{gauge orbits} the function --- and hence $L_{\mathrm{sep}}$ --- is exactly
constant, so their tangent directions lie in the kernel of the factored Hessian:
\begin{equation}\label{eq:nulldim}
\dim\ker\bigl(\Hess_\theta L_{\mathrm{sep}}\bigr)\ \ge\ r^2+r'^2 \quad\text{per layer.}
\end{equation}
These $r^2+r'^2$ flat directions exist purely because of the
over‑parametrization; the combined coordinates $M$ have no such redundancy.
Gradient descent conserves the imbalance $\|\Wq\|_F^2-\|\Wk\|_F^2$ (and likewise
for OV) \cite{DuHuLee2018}, so it lives on a fixed --- the \emph{balanced} ---
slice of each orbit; this is the canonical point at which any fair comparison of
the two parametrizations must be made.

\section{Main result}

Collecting \eqref{eq:hess}, Lemma~\ref{lem:bilin}, \eqref{eq:Rspec} and
\eqref{eq:nulldim}:

\begin{theorem}[Factoring only adds curvature pathologies]\label{thm:main}
Let $P=J_\varphi^{\!\top}(\Hess_M\widetilde L)\,J_\varphi$ and $R$ be as in
\eqref{eq:hess}, so that $\Hess_\theta L_{\mathrm{sep}}=P+R$.
\begin{enumerate}
\item[\textup{(a)}] \textbf{(PSD pullback with extra flats.)}
If $\Hess_M\widetilde L\succeq 0$ then $P\succeq 0$. Moreover
$\ker P\supseteq$ the gauge tangent space of \eqref{eq:gauge}, of dimension
$\ge r^2+r'^2$ per layer.
\item[\textup{(b)}] \textbf{(Indefinite residual off-stationarity.)}
$R$ has spectrum \eqref{eq:Rspec}; it is symmetric about $0$, and
\[
\lambda_{\min}(R)\;=\;-\max\!\Bigl\{\sigma_{\max}\!\bigl(\nabla_{\Wqk}\widetilde L\bigr),\,
                                      \sigma_{\max}\!\bigl(\nabla_{\Wov}\widetilde L\bigr)\Bigr\}\;<\;0
\quad\text{whenever }\nabla_M\widetilde L\neq 0 .
\]
\item[\textup{(c)}] \textbf{(Comparison.)} The factored Hessian differs from the
PSD pullback of the combined Hessian by exactly the indefinite operator $R$:
\[
\Hess_\theta L_{\mathrm{sep}} - P \;=\; R,\qquad
\lambda_{\min}(R)<0\ \text{off every stationary point.}
\]
At a stationary point ($\nabla_M\widetilde L=0$) we have $R=0$, so
$\Hess_\theta L_{\mathrm{sep}}=P$ agrees with $\Hess_M\widetilde L$ on the
non‑gauge subspace and differs only by the flat gauge null space.
\end{enumerate}
\end{theorem}

\begin{remark}[What the theorem says in words]
The combined coordinates see the bare circuit curvature $\Hess_M\widetilde L$.
The factored coordinates see (i)~a \emph{stretched, rank‑deficient} version of it
(the Gauss--Newton pullback $P$, which also adds $r^2{+}r'^2$ flat valleys), plus
(ii)~an \emph{extra indefinite term} $R$ that plants a saddle of depth
$\sigma_{\max}(\nabla_M\widetilde L)$ at every non‑optimal point. Factoring can
therefore only \emph{remove convexity, never add it}, and is strictly worse
wherever the gradient is nonzero. This is the induction‑circuit instance of the
classical matrix‑factorization landscape: products of parameters create saddles
at non‑optimal points
\cite{BaldiHornik1989,Kawaguchi2016,GeLeeMa2016,SaxeMcCGanguli2014}.
\end{remark}

\section{Intuition}\label{sec:intuition}

\paragraph{Why combined is gentle.}
In the combined coordinates the score is \emph{linear} in $\Wqk$ and the value
output is \emph{linear} in $\Wov$ (only the $\softmax$ in between is nonlinear).
Each circuit is a single coordinate with no internal degrees of freedom, so the
only curvature is the irreducible curvature of the task itself. Empirically this
makes the combined Hessian almost isotropic --- a nearly spherical bowl that
gradient descent can roll down with a single, large step size.

\paragraph{Why factoring hurts.}
Writing a circuit as a \emph{product} $\Wq^{\!\top}\Wk$ introduces multiplicative
coupling between the two factors, the hallmark of deep‑linear optimization. Two
things follow immediately and are visible in \eqref{eq:gauge} and \eqref{eq:Rspec}:
\begin{itemize}
\item \textbf{Flat valleys (gauge).} Scaling $\Wq$ up and $\Wk$ down by a matching
factor leaves the product --- and the loss --- unchanged. These are the
$r^2{+}r'^2$ zero modes of \eqref{eq:nulldim}: long, flat troughs that
ill‑condition the problem without changing the function.
\item \textbf{Saddles (residual).} Away from the optimum the reduced gradient
$G=\nabla_{\Wqk}\widetilde L$ is nonzero, and the move ``increase $\Wq$ along its
top singular direction while decreasing $\Wk$'' is a direction of
\emph{negative} curvature of depth $\sigma_{\max}(G)$ (Lemma~\ref{lem:bilin}).
The factored surface bends \emph{down} in directions where the combined surface
is flat or up.
\end{itemize}
So GD in the factored coordinates must thread between flat gauge troughs and
strict saddles, and must use a tiny step size because the same multiplicative
coupling also creates very sharp directions; GD in the combined coordinates sees
none of this.

\section{Empirical confirmation}\label{sec:empirics}

We took the trained combined checkpoints, produced the \emph{balanced} (SVD)
factorization that realizes the identical function, and computed the Hessian of
the suffix‑$100$ cross‑entropy in both coordinate systems via Hessian--vector
products and power iteration (window of $128$ test sequences). The combined
($400{,}420$ parameters) and factored ($739{,}220$ parameters --- the extra
$338{,}800$ are exactly the gauge redundancy) nets agree to floating‑point on the
loss. Table~\ref{tab:hess} reports, at the trained solution and at an
off‑solution point (all weights scaled by $0.7$, where $\nabla_M\widetilde L\neq0$).

\begin{table}[h]
\centering
\small
\begin{tabular}{llrrrr}
\toprule
$\alpha$ & point & param. & $\lambda_{\max}$ (sharpness) & $\lambda_{\min}$ & $\kappa\approx\lambda_{\max}/\lambda^{+}$ \\
\midrule
$0.9$ & solution      & combined & $3.4\times10^{-5}$ & $+2.3\times10^{-5}$ (PSD) & $\mathbf{\approx 1.3}$\\
$0.9$ & solution      & factored & $\ge 1.21$          & $+7\times10^{-5}$         & $\mathbf{\approx 1.4\times10^{4}}$\\
$0.9$ & off‑solution  & combined & $2.1\times10^{-5}$ & $+3.8\times10^{-6}$ (PSD)  & $\approx 3.6$\\
$0.9$ & off‑solution  & factored & ---                & $\mathbf{-0.74}$ (saddle)  & non‑convex\\
\midrule
$0.5$ & solution      & combined & $6.8\times10^{-5}$ & $+5.5\times10^{-5}$ (PSD) & $\mathbf{\approx 1.2}$\\
$0.5$ & solution      & factored & $\ge 3.11$          & $+1\times10^{-4}$         & $\mathbf{\approx 3.0\times10^{4}}$\\
\bottomrule
\end{tabular}
\caption{Hessian geometry of the same function in the two parametrizations.
$\lambda^{+}$ is a gauge‑robust lower curvature scale (the $10$th percentile of
random Rayleigh quotients), used because the factored Hessian has many exact
gauge zeros \eqref{eq:nulldim}. The combined Hessian is near‑isotropic
($\kappa\approx1$) and stays PSD off the optimum; the factored Hessian is
$\sim\!10^4$ worse conditioned and develops a negative‑curvature saddle
($\lambda_{\min}=-0.74$) off the optimum, exactly as Theorem~\ref{thm:main}(b)
predicts.}
\label{tab:hess}
\end{table}

The numbers match the theory point for point: (i)~at the solution
$\nabla_M\widetilde L\approx0$, so $R\approx0$ and both forms are PSD --- but the
factored form is already $\sim\!10^4\times$ sharper (the pullback $P$ stretches
curvature and adds gauge flats); (ii)~off the solution $\nabla_M\widetilde L\neq0$,
and the factored form acquires the predicted negative eigenvalue while the
combined form remains positive.

\section{Caveats}

\begin{enumerate}
\item \textbf{Gauge.} ``Condition number'' for the factored form excludes the
exact gauge zeros of \eqref{eq:nulldim}; the comparison is made at the balanced
point where gradient descent actually sits \cite{DuHuLee2018}.
\item \textbf{``More convex'' $\neq$ ``always trains better.''} Over‑parametrized
factored forms can also \emph{implicitly accelerate} and bias toward low rank
\cite{AroraCohenHazan2018}, and the injected saddles are \emph{strict}, hence
escapable by perturbed/stochastic GD \cite{GeHuangJinYuan2015}. The rigorous claim
here is about \emph{conditioning and saddle‑proneness}, which is exactly what the
folder‑7 combined parametrization buys by removing the within‑head rotation gauge.
\end{enumerate}

\begin{thebibliography}{9}
\bibitem{BaldiHornik1989} P.~Baldi and K.~Hornik.
\emph{Neural networks and principal component analysis: learning from examples
without local minima.} Neural Networks, 1989.
\bibitem{SaxeMcCGanguli2014} A.~Saxe, J.~McClelland, S.~Ganguli.
\emph{Exact solutions to the nonlinear dynamics of learning in deep linear neural
networks.} ICLR, 2014.
\bibitem{Kawaguchi2016} K.~Kawaguchi.
\emph{Deep learning without poor local minima.} NeurIPS, 2016.
\bibitem{GeHuangJinYuan2015} R.~Ge, F.~Huang, C.~Jin, Y.~Yuan.
\emph{Escaping from saddle points --- online stochastic gradient for tensor
decomposition.} COLT, 2015.
\bibitem{GeLeeMa2016} R.~Ge, J.~D.~Lee, T.~Ma.
\emph{Matrix completion has no spurious local minimum.} NeurIPS, 2016.
\bibitem{DuHuLee2018} S.~S.~Du, W.~Hu, J.~D.~Lee.
\emph{Algorithmic regularization in learning deep homogeneous models: layers are
automatically balanced.} NeurIPS, 2018.
\bibitem{AroraCohenHazan2018} S.~Arora, N.~Cohen, E.~Hazan.
\emph{On the optimization of deep networks: implicit acceleration by
overparameterization.} ICML, 2018.
\bibitem{KarimiNutiniSchmidt2016} H.~Karimi, J.~Nutini, M.~Schmidt.
\emph{Linear convergence of gradient and proximal‑gradient methods under the
Polyak--{\L}ojasiewicz condition.} ECML‑PKDD, 2016.
\bibitem{Yao2020} Z.~Yao, A.~Gholami, K.~Keutzer, M.~Mahoney.
\emph{PyHessian: neural networks through the lens of the Hessian.} IEEE BigData,
2020.
\end{thebibliography}

\end{document}
